% chktex-file 8
% chktex-file 12
% chktex-file 18
\beginsong{Čarodějnice z Amesbury}[by={Asonance}]

\beginverse{}
Zuzana \ch{Dmi}{}{}{}byla dívka, \ch{C}{}{}{}která žila \ch{Dmi}{}{}{}v Amesbury,
s jasnýma \ch{F}{}{}{}očima a \ch{C}{}{}{}řečmi pánům \ch{Dmi}{}{}{}navzdory,
souse\ch{F}{}{}{}dé o ní \ch{C}{}{}{}říkali, že \ch{Dmi}{}{}{}temná kouzla \ch{Ami}{}{}{}zná
a že \ch{B{\flt}}{}{}{}se lidem \ch{Ami}{}{}{}vyhýbá a s \ch{B{\flt}}{}{}{}ďáblem \ch{C}{}{}{}pletky \ch{Dmi}{}{}{}má.
\endverse{}

\beginverse{}
Onoho léta náhle mor dobytek zachvátil
a pověrčivý lid se na pastora obrátil,
že znají tu moc nečistou, jež krávy zabíjí,
a odkud ta moc vychází, to každý dobře ví.
\endverse{}

\beginverse{}
Tak Zuzanu hned před tribunál předvést nechali,
a když ji vedli městem, všichni kolem volali
"Už konec je s tvým řáděním, už nám neuškodíš,
teď na své cestě poslední do pekla poletíš!"
\endverse{}

\beginverse{}
Dosvědčil jeden sedlák, že zná její umění,
ďábelským kouzlem prý se v netopýra promění
a v noci nad krajinou létává pod černou oblohou,
sedlákům krávy zabíjí tou mocí čarovnou.
\endverse{}

\beginverse{}
Jiný zas na kříž přísahal, že její kouzla zná,
v noci se v černou kočku mění dívka líbezná,
je třeba jednou provždy ukončit ďábelské řádění,
a všichni křičeli jako posedlí:"Na šibenici s ní!"
\endverse{}

\beginverse{}
Spektrální důkazy pečlivě byly zváženy,
pak z tribunálu povstal starý soudce vážený:
"Je přece v knize psáno: nenecháš čarodějnici žít
a před ďáblovým učením budeš se na pozoru mít!"
\endverse{}

\beginverse{}
Zuzana stála krásná s hlavou hrdě vztyčenou
a její slova zněla klenbou s tichou ozvěnou:
"Pohrdám vámi, neznáte nic nežli samou lež a klam,
pro tvrdost vašich srdcí jen, jen pro ni umírám!"
\endverse{}

\beginverse{}
Tak vzali Zuzanu na kopec pod šibenici
a všude kolem ní se sběhly davy běsnící,
a ona stála bezbranná, však s hlavou vztyčenou,
zemřela tiše samotná pod letní oblohou..
\endverse{}

\endsong{}
